\scanpage{50}\mspage{3\deleted{\unsure{8|6}}7}

\uline{Lemme} 2.9. $I$ ens. ordonné quelconque, $C,W$ comme dans Th.\ 2.1, $1^\circ$ et $3^\circ$, $\deleted{\unsure{v}}u_\bullet:X_\bullet\longrightarrow Y_\bullet$ dans $C_I\cap W_{I^{\unsure{o|\circ}}}$. Soit $P$ l'ens des $i_0\in I$ tels que $u_\bullet\mid I_{/i_0}$ soit dans $(TC)_{I_{/i_0}}$. \deleted{Alors} Si $i_0\in I$ est tel que $I_{<i_0}\unsure{\in|\subset}P$, alors $i_0\in P$.

En effet, \deleted{\unsure{on aura} $u_\bullet\mid$} il reste à prouver que pour \uline{tout} ouvert $K$ de $J=I_{/i_0}$, \inserted{$K\ne J$,} on a $u_{K,J}\in TC$. Mais \deleted{\unsure{on a} $K\subset J'=I_{<i_0}$, et il suffit de voir que $u_{K,J'}\in TC$, $u_{J',J}\in W$.} Pour ceci, il suffit de voir que $u_K,u_J\in\deleted{\unsure{TC|W}}\unsure{W|TC}$ (comme on vient de le voir). Or $u_J\simeq u_{i_0}\in W$ par hypothèse.

D'autre part $\deleted{\unsure{u_K\in}}u_\bullet\mid K\in(TC)_K$ \deleted{par hyp. sur} puisque $K\subset I_{<i_0}\subset P$, donc par le Th.\ 1, $1^\circ$, on conclut que $u_K\in TC$.

Si maintenant on suppose $I$ artinien, il résulte de 3.9. que $P=I$, donc $u_\bullet\in TC_I$, qed.

Cela achève la démonstration du Th.\ 2.1, et aussi du cor.\ 2.5, vu que pour la \inserted{\unsure{dém.}} \inserted{du} lemme 2.9. également, les hypothèses infinies sur $C$, $C\cap W$ sont inutiles si $\operatorname{Top}(I)$ est noethérien, donc ses ouverts \deleted{\unsure{tt}} \deleted{\unsure{étant}} sont \unsure{aussi} noethériens.

\uline{Remarque}\deleted{\unsure{s}} $1)$ Dans 2.1 $2^\circ$, l'hypothèse $I$ artinien n'est pas superflue, comme on voit en prenant $I=\mathbb N^{\mathrm{opp}}\simeq\deleted{\unsure{\mathbb Z}}\mathbb Z^{\le0}$,

% archive sheet number: 49

